\documentclass[11pt]{article}
\usepackage[a4paper,margin=28mm]{geometry}
\usepackage{amsmath,amssymb,amsthm,hyperref}
\newtheorem{theorem}{Theorem}
\newtheorem{proposition}[theorem]{Proposition}
\newtheorem{lemma}[theorem]{Lemma}
\newcommand{\Z}{\mathbb Z}
\newcommand{\Q}{\mathbb Q}
\newcommand{\hgt}{\operatorname{h}}
\title{Exponential growth from two integer labels}
\author{}
\date{9 November 2026}
\begin{document}
\maketitle
\begin{abstract}
An indefinite norm equation produces exponential growth using two
integer labels and one rational root of a monic polynomial. We give an
elementary descent proof and a precise limitation on applying the
construction to rational approximations of the integer labels. No
reduction in the number of variables of a defining polynomial is
asserted.
\end{abstract}

For an integer $a\ge3$, define
\[
 U_0(a)=0,\qquad U_1(a)=1,\qquad
 U_{j+1}(a)=aU_j(a)-U_{j-1}(a)\quad(j\ge1).
\]

\begin{lemma}\label{lem:descent}
If $v,w\in\Z$ satisfy
\[
                       w^2-avw+v^2=1,
\]
then $|v|=U_j(a)$ for some $j\ge0$.
For $j\ge1$ one has
\[
 U_j(a)\equiv j\pmod{a-2},\qquad
 U_j(a)\ge(a-1)^{j-1}.
\]
\end{lemma}
\begin{proof}
The congruence follows by induction because $a\equiv2\pmod{a-2}$.
The sequence is strictly increasing from $U_1$ onwards, and
$U_{j+1}=(a-1)U_j+(U_j-U_{j-1})\ge(a-1)U_j$ for $j\ge1$.
This proves the lower bound.

If $vw<0$, the left-hand side of the norm equation is larger than
one. If $vw=0$, the other coordinate is $\pm1$. Otherwise the two
coordinates have the same sign, and we can change both signs to
make them positive. The equation is invariant under interchange
of the coordinates and under
\[
                         (v,w)\longmapsto(v,av-w).
\]
Order the coordinates as $1\le v\le w$. If $v>1$, then
\[
                  0<av-w=\frac{v^2-1}{w}<v.
\]
Thus interchange followed by this operation decreases the maximum
coordinate. It terminates with $v=1$, where $w(w-a)=0$.
Reversing these operations gives the consecutive pairs
$(U_j(a),U_{j+1}(a))$. Every coordinate of the original pair is
therefore one of the $U_j(a)$.
\end{proof}

\begin{theorem}\label{thm:growth}
Let $a,u\in\Z$, with $a\ge7$ and $u\ne0$, and put
$v=(a-2)u+1$. If a rational number $w$ satisfies
\begin{equation}\label{eq:monic}
                      w^2-avw+v^2-1=0,
\end{equation}
then $w\in\Z$, and
\[
          |v|\ge(a-1)^{a-4},\qquad
          |u|\ge\frac{(a-1)^{a-4}-1}{a-2}.
\]
Conversely, for every integer $a\ge7$ there exist integers $u\ne0$
and $w$ satisfying \eqref{eq:monic}.
\end{theorem}
\begin{proof}
The polynomial in \eqref{eq:monic} is monic with integer
coefficients, so its rational root is integral. The conditions
$a\ge7$ and $u\ne0$ exclude $v=0,1,-1$: the first and third would
require $a-2$ to divide $1$ or $2$, and the second requires $u=0$.
By Lemma~\ref{lem:descent}, $|v|=U_j(a)$ for an integer $j\ge2$.
Since $v\equiv1\pmod{a-2}$, one has
$j\equiv1$ or $-1\pmod{a-2}$. Hence $j\ge a-3$ and
$|v|\ge(a-1)^{a-4}$. The bound for $u$ follows from
$|v|\le(a-2)|u|+1$.

For existence, take
\[
 v=U_{a-1}(a),\qquad w=U_a(a),\qquad
 u=\frac{U_{a-1}(a)-1}{a-2}.
\]
The congruence in Lemma~\ref{lem:descent} makes $u$ an integer,
and $u>0$. The recurrence preserves
$U_{j+1}^2-aU_jU_{j+1}+U_j^2=1$, as follows by induction from
$(U_0,U_1)=(0,1)$.
\end{proof}

The theorem reduces the number of integer labels needed for this
growth mechanism from three to two. The root $w$ is a rational
variable rather than an additional integer label. This distinction
is useful only when the coefficients in \eqref{eq:monic} are the
actual integers $a,u$. In a construction using rational
approximations of those labels, integrality of $w$ is not available
until exact equality of the coefficients has been proved.

\begin{proposition}\label{prop:approximation}
Fix $a\ge7$ and any nonzero integral solution $(v_0,w_0)$ of the
norm equation with $v_0\equiv1\pmod{a-2}$, and put
$u_0=(v_0-1)/(a-2)$. Outside a finite set of primes $p$, the
following assertion holds for every positive integer $k$.
There are rational solutions $(u',w')$ with
\[
 v'=(a-2)u'+1,\qquad (w')^2-av'w'+(v')^2=1,
 \qquad u'\equiv u_0,\quad w'\equiv w_0\pmod{p^k},
\]
for which the logarithmic rational height $\hgt(w')$ is unbounded.
\end{proposition}
\begin{proof}
Use the rational parametrization through $(v,w)=(0,1)$:
\[
 v(s)=\frac{a-2s}{s^2-as+1},\qquad
 w(s)=\frac{1-s^2}{s^2-as+1}.
\]
It follows by substituting the line $w=1+sv$ into the norm
equation. Since $v_0\ne0$, the parameter
$s_0=(w_0-1)/v_0$ is rational. Its denominator
$s_0^2-as_0+1$ is nonzero: otherwise $s_0=a/2$ would also be
necessary, contrary to $a^2\ne4$.
Exclude primes dividing the denominators of $s_0$,
$s_0^2-as_0+1$, their numerators when necessary for the latter
to be a unit, or $a-2$. Both rational functions and
$u(s)=(v(s)-1)/(a-2)$ are then regular on a $p$-adic
neighborhood of $s_0$ and preserve congruences modulo $p^k$.

Take $s=s_0+p^k/m$ with positive integers $m$ prime to $p$.
These give infinitely many distinct rational parameters with the
required congruences. At most two parameters are poles. The
rational function $w(s)$ is nonconstant, and every value has at
most two inverse images. Its values on this sequence therefore
form an infinite set. Rational numbers of bounded logarithmic
height form a finite set, so their heights are unbounded.
\end{proof}

\section*{The remaining arithmetic requirement}
In the existing comparison construction, exact integral point
indices are recovered from rational formal-parameter ratios by
divisibility of an integer residual, whose ordinary absolute value
has a polynomial bound in those indices. Replacing a Lucas label
by the monic root in Theorem~\ref{thm:growth} removes this bound
unless the height or denominator of the rational root is separately
controlled. Proposition~\ref{prop:approximation} shows that any
fixed amount of local agreement alone cannot provide such control.
It does not rule out a stronger argument using the global geometry
of the retained elliptic points. A successful application needs an
additional global bound, or another means of proving the exact
integral norm equation before invoking monicity.

This issue persists when the contact-height estimate is strengthened
from one fixed positive power to another. To make the limitation
precise, suppose the available bounds have the shape
\[
 h(t)\le C M(t)^\alpha,\qquad
 B^2\le C'|J|^{\alpha/k}+C'',\qquad
 |J|\le2B\exp h(\beta),
\]
where $\alpha>0$ and $k>\alpha/2$ are fixed. For sufficiently large
$B$, taking logarithms yields only
\[
 \left(2-\frac\alpha k\right)\log B
 \le O(1)+\frac\alpha k h(\beta).
\]
The scalar conclusion $h(\beta)\le C_0c^\epsilon$, with any fixed
$\epsilon>0$, and the elementary inequality $B\ge c-1$ are
compatible with this bound for arbitrarily large $c$. Decreasing
$\alpha$ or increasing a fixed precision does not, by itself,
replace the exponential-growth step.

Likewise, adjoining scalar scales that merely bound the rational
root height by $h(w)\le C_1M^\alpha$ does not close the residual
argument. If the residual numerator is bounded by
$B^r\exp(s h(w))$ for fixed $r,s>0$, divisibility gives
$k\log M\le r\log B+s h(w)$, which is consistent with
$h(w)$ growing as a positive power of $M$. These observations do
not exclude other uses of a stronger contact estimate. They identify
the missing conclusion in this particular replacement: a bound such
as $h(w)=O(\log B)$, or exact integrality established before comparing
the residual with its local divisors.
\end{document}
